\section{System Architecture}

The MediPilot system is designed around a strictly separated, parallel track architecture. This separation of concerns ensures that the machine learning models operate within tight bounds and that the core clinical safety rules can never be bypassed by unpredictable AI outputs.

\subsection{The Fail-Closed Philosophy}
In clinical artificial intelligence, a ``Fail-Closed'' architecture dictates that in the event of missing data, out-of-distribution inputs, or model uncertainty, the system defaults to the safest deterministic clinical rule rather than generating a potentially hallucinated or miscalibrated output. 

MediPilot achieves this by placing hard clinical rules (such as red flags and multi-vital distress floors) \textit{outside} of the model path. If the model fails or is bypassed, the rule card governs the patient's triage band.

\subsection{Architectural Tracks}
The system is divided into three parallel tracks:
\begin{itemize}
    \item \textbf{Track A (Perception Layer):} Owns everything between the patient/nurse's voice and a clean, structured record. It handles speech recognition, age-aware questioning, and grammar-constrained Large Language Model (LLM) extraction.
    \item \textbf{Track B (Nurse Frontend):} Manages direct vital sign entry and physical kiosk interactions. (Not detailed in this document).
    \item \textbf{Track C (Risk Engine \& Backend):} Receives structured records from Track A and vitals from Track B to orchestrate triage band assignment (Red/Yellow/Green), recheck scheduling, and hospital surge control using a combination of deterministic rules and a Gradient Boosted Decision Tree (GBDT).
\end{itemize}
